% Part: turing-machines
% Chapter: machines-computations
% Section: disciplined-machines

\documentclass[../../../include/open-logic-section]{subfiles}

\begin{document}

\olfileid[fr]{tur}{mac}{dis}
\olsection{Machines disciplinées}

\begin{explain}
Dans la \olref[hal]{sec}, nous avons considéré des machines de
Turing dotées d'un unique état d'arrêt désigné~$h$ : si elles
s'arrêtent, c'est nécessairement dans cet état. Elles sont
en ce sens plus « disciplinées » que les machines de Turing
que nous autorisons en général. On pourrait imposer d'autres
restrictions à leur comportement. Par exemple, nous n'avons
pas interdit à une machine d'effacer le marqueur de début du
ruban dans la case~$0$, ni de tenter de déplacer sa tête vers
la gauche depuis cette case. Notre définition prévoit que
la tête reste alors dans la case~$0$; dans d'autres définitions,
la machine s'arrête. On souhaite parfois aussi pouvoir supposer
que, si une machine s'arrête, sa tête lit la case~$1$.
\end{explain}

\begin{defn}\ollabel{defn:disciplined}
Une machine de Turing~$M$ est \emph{disciplinée} si et
seulement si :
\begin{enumerate}
    \item elle possède un unique état d'arrêt désigné~$h$;
    \item si elle s'arrête, sa tête lit alors la case~$1$;
    \item elle n'efface jamais le symbole $\TMendtape$ de
      la case~$0$;
    \item elle ne tente jamais de déplacer sa tête vers la
      gauche depuis la case~$0$.
\end{enumerate}
\end{defn}

\begin{explain}
Nous avons déjà expliqué comment transformer une machine de
Turing en une machine de même comportement quant à l'arrêt
et à la sortie, munie d'un état
d'arrêt désigné. On ajoute un nouvel état~$h$ puis, pour toute
paire $\tuple{q,\sigma}$ d'un état de la machine initiale et
d'un symbole sur laquelle sa fonction~$\delta$ n'est pas
définie, on ajoute l'instruction
$\delta(q,\sigma)=\tuple{h,\sigma,N}$. Aucune instruction
ne part du nouvel état~$h$.\footnote{Précision éditoriale : les transitions sont ajoutées seulement pour les états anciens ; en ajouter au nouvel état $h$ empêcherait celui-ci d'être un état d'arrêt.}

On peut aussi, bien que la preuve soit fastidieuse, transformer
toute machine de Turing~$M$ en une machine disciplinée~$M'$
qui s'arrête sur les mêmes entrées et produit les mêmes sorties.
Par exemple, si la machine s'arrête alors que sa tête ne lit
pas la case~$1$, on peut ajouter des instructions qui ramènent
la tête vers un marqueur réservé à la case~$0$, puis la
déplacent d'une case vers la droite avant l'arrêt. Il faut
également faire en sorte que ce marqueur soit préservé et
reconnaissable pendant la simulation.\footnote{Précision éditoriale : le marqueur source peut être écrasé ou apparaître ailleurs. Le retour vers la case $0$ suppose donc un marqueur réservé et préservé par la simulation ; il ne suffit pas de chercher une occurrence quelconque de l'ancien symbole.} Nous vous laissons
réfléchir au traitement des autres conditions.
\end{explain}

\begin{ex}
    La machine d'addition de l'\olref[una]{ex:adder} est
    transformée ici en machine disciplinée. Le résultat apparaît
    dans la \olref{fig:adder-disc}.\footnote{Correction éditoriale : comme dans le premier diagramme d'addition, la transition source de $q_0$ lisant un trait mène à $q_1$ ; nous la corrigeons en une boucle sur $q_0$.}
    \begin{figure}\[
\begin{tikzpicture}[->,>=stealth',shorten >=1pt,auto,node distance=2.8cm,
                    semithick]
  \tikzstyle{every state}=[fill=none,draw=black,text=black]

  \node[initial,state]         (A)              {$q_0$};
  \node[state]         (B) [right of=A] {$q_1$};
  \node[state]         (C) [below right of=B] {$q_2$};
  \node[state]         (D) [below left of=C] {$q_3$};
  \node[state]         (H) [left of=D] {$h$};

  \path (A) edge node {\TMtrans{\TMblank}{\TMstroke}{\TMstay}} (B)
           edge [loop below] node {\TMtrans{\TMstroke}{\TMstroke}{\TMright}} (A)
        (B) edge [loop below] node {\TMtrans{\TMstroke}{\TMstroke}{\TMright}} (B)
            edge node {\TMtrans{\TMblank}{\TMblank}{\TMleft}} (C)
        (C) edge node {\TMtrans{\TMstroke}{\TMblank}{\TMleft}} (D)
        (D) edge [loop above] node {\TMtrans{\TMstroke}{\TMstroke}{\TMleft}} (D)
        edge node {\TMtrans{\TMendtape}{\TMendtape}{\TMright}} (H);
\end{tikzpicture}
\]\caption{Machine d'addition disciplinée}
\ollabel{fig:adder-disc}
\end{figure}
\end{ex}

\begin{prop}\ollabel{prop:disciplined}
Pour toute machine de Turing~$M$, il existe une machine
de Turing disciplinée~$M'$ telle que, si $M$ s'arrête avec
la sortie~$O$, alors $M'$ s'arrête avec la sortie~$O$,
et que, si $M$ ne s'arrête pas, $M'$ ne s'arrête pas non plus.
En particulier, toute fonction
$f\colon\Nat^n \to \Nat$ calculable par une machine de
Turing est aussi calculable par une machine disciplinée.
\end{prop}

\begin{prob}\label{tur:mac:dis:prob:disc-succ}
Donnez une machine disciplinée qui calcule $f(x)=x+1$.
\end{prob}

\begin{prob}\label{tur:mac:dis:prob:copier}
Trouvez une machine disciplinée qui, sur l'entrée
$\TMstroke^n$, produit la sortie
$\TMstroke^n \concat \TMblank \concat \TMstroke^n$.
\end{prob}

\end{document}
